% Options for packages loaded elsewhere
\PassOptionsToPackage{unicode}{hyperref}
\PassOptionsToPackage{hyphens}{url}
%
\documentclass[
  11pt,
]{article}
\usepackage{lmodern}
\usepackage{amssymb,amsmath}
\usepackage{ifxetex,ifluatex}
\ifnum 0\ifxetex 1\fi\ifluatex 1\fi=0 % if pdftex
  \usepackage[T1]{fontenc}
  \usepackage[utf8]{inputenc}
  \usepackage{textcomp} % provide euro and other symbols
\else % if luatex or xetex
  \usepackage{unicode-math}
  \defaultfontfeatures{Scale=MatchLowercase}
  \defaultfontfeatures[\rmfamily]{Ligatures=TeX,Scale=1}
\fi
% Use upquote if available, for straight quotes in verbatim environments
\IfFileExists{upquote.sty}{\usepackage{upquote}}{}
\IfFileExists{microtype.sty}{% use microtype if available
  \usepackage[]{microtype}
  \UseMicrotypeSet[protrusion]{basicmath} % disable protrusion for tt fonts
}{}
\makeatletter
\@ifundefined{KOMAClassName}{% if non-KOMA class
  \IfFileExists{parskip.sty}{%
    \usepackage{parskip}
  }{% else
    \setlength{\parindent}{0pt}
    \setlength{\parskip}{6pt plus 2pt minus 1pt}}
}{% if KOMA class
  \KOMAoptions{parskip=half}}
\makeatother
\usepackage{xcolor}
\IfFileExists{xurl.sty}{\usepackage{xurl}}{} % add URL line breaks if available
\IfFileExists{bookmark.sty}{\usepackage{bookmark}}{\usepackage{hyperref}}
\hypersetup{
  pdftitle={The Analytic Atlas of the UNNS Substrate},
  pdfauthor={UNNS Research Collective},
  hidelinks,
  pdfcreator={LaTeX via pandoc}}
\urlstyle{same} % disable monospaced font for URLs
\usepackage[margin=25mm]{geometry}
\usepackage{longtable,booktabs}
% Correct order of tables after \paragraph or \subparagraph
\usepackage{etoolbox}
\makeatletter
\patchcmd\longtable{\par}{\if@noskipsec\mbox{}\fi\par}{}{}
\makeatother
% Allow footnotes in longtable head/foot
\IfFileExists{footnotehyper.sty}{\usepackage{footnotehyper}}{\usepackage{footnote}}
\makesavenoteenv{longtable}
\setlength{\emergencystretch}{3em} % prevent overfull lines
\providecommand{\tightlist}{%
  \setlength{\itemsep}{0pt}\setlength{\parskip}{0pt}}
\setcounter{secnumdepth}{-\maxdimen} % remove section numbering

\title{The Analytic Atlas of the UNNS Substrate}
\author{UNNS Research Collective}
\date{27 November 2026}

\begin{document}
\maketitle

{
\setcounter{tocdepth}{1}
\tableofcontents
}
\hypertarget{how-to-read-the-atlas}{%
\section{How to read the Atlas}\label{how-to-read-the-atlas}}

The Atlas organizes the results of Analytic Structures of the UNNS
Substrate by mathematical phenomenon. It covers ANS-I through v015,
ANS-II through v002 and ANS-II-ND through v003. It states nothing new:
each mathematical claim is a short account of catalogued results, named
by identifier. The exact statements, with all hypotheses, are in the
Theorem Compendium, and they govern wherever a summary here is shorter.

Each entry gives the mathematical object, the established structure,
what that structure says about generation, further structure and the
open boundary. The sections headed \textbf{Reading across} are
interpretation: they point to parallels and are not theorems. The
tables, the connections between entries, the review boundary and the
proof sources at the end of each entry are generated from the catalogue.

Status words: \emph{proved} means a written argument exists; review
codes say who has checked it (V UNNS Research Collective review, VQ UNNS
Research Collective review with a qualification, A proposal checked
internally, P no review verdict, B baseline note, N not a proved
result). No literature novelty review has been completed.

\textbf{Notation.} A letter \(N\ge1\) multiplies by \(r_N=(N+1)^2/N\)
and is admissible at \(x\) when \(N\mid x\). For a word \(w\):
multiplier \(Q(w)\), minimal source \(K(w)\), least endpoint
\(E(w)=K(w)Q(w)\). Further, \(G_k(s)=\sum_{|w|=k}E(w)^{-s}\),
\(\mathcal G(s,t)=\sum_kt^kG_k(s)\), \(B(\sigma)=\sum_Nr_N^{-\sigma}\),
and \(\sigma_c=1.5463810\ldots\) solves \(B(\sigma_c)=1\).

\hypertarget{four-governing-mechanisms}{%
\section{Four governing mechanisms}\label{four-governing-mechanisms}}

The foundational audit of 27 November 2026 singled out four results on
which much of the theory turns.

\begin{longtable}[]{@{}llll@{}}
\toprule
\begin{minipage}[b]{0.23\columnwidth}\raggedright
Mechanism\strut
\end{minipage} & \begin{minipage}[b]{0.13\columnwidth}\raggedright
Result\strut
\end{minipage} & \begin{minipage}[b]{0.41\columnwidth}\raggedright
What it does\strut
\end{minipage} & \begin{minipage}[b]{0.11\columnwidth}\raggedright
Entry\strut
\end{minipage}\tabularnewline
\midrule
\endhead
\begin{minipage}[t]{0.23\columnwidth}\raggedright
Exact admissible composition\strut
\end{minipage} & \begin{minipage}[t]{0.13\columnwidth}\raggedright
ANS-T03\strut
\end{minipage} & \begin{minipage}[t]{0.41\columnwidth}\raggedright
Turns histories into arithmetic: every word has an exact set of sources
and endpoints\strut
\end{minipage} & \begin{minipage}[t]{0.11\columnwidth}\raggedright
ATL-01\strut
\end{minipage}\tabularnewline
\begin{minipage}[t]{0.23\columnwidth}\raggedright
Finite-resource saturation\strut
\end{minipage} & \begin{minipage}[t]{0.13\columnwidth}\raggedright
ANS-T15\strut
\end{minipage} & \begin{minipage}[t]{0.41\columnwidth}\raggedright
Each finite truncation \(1\le N\le M\) can be prepared so that random
generation continues for ever with positive probability\strut
\end{minipage} & \begin{minipage}[t]{0.11\columnwidth}\raggedright
ATL-02\strut
\end{minipage}\tabularnewline
\begin{minipage}[t]{0.23\columnwidth}\raggedright
Unbounded-resource obstruction\strut
\end{minipage} & \begin{minipage}[t]{0.13\columnwidth}\raggedright
ANS-T17\strut
\end{minipage} & \begin{minipage}[t]{0.41\columnwidth}\raggedright
With all letters available, a random infinite word is almost surely
inadmissible from any integer\strut
\end{minipage} & \begin{minipage}[t]{0.11\columnwidth}\raggedright
ATL-02\strut
\end{minipage}\tabularnewline
\begin{minipage}[t]{0.23\columnwidth}\raggedright
Quantitative collective suppression\strut
\end{minipage} & \begin{minipage}[t]{0.13\columnwidth}\raggedright
ANS-T23a\strut
\end{minipage} & \begin{minipage}[t]{0.41\columnwidth}\raggedright
Measures the cost of that obstruction at depth \(k\)\strut
\end{minipage} & \begin{minipage}[t]{0.11\columnwidth}\raggedright
ATL-04\strut
\end{minipage}\tabularnewline
\bottomrule
\end{longtable}

\hypertarget{the-shape-of-the-territory}{%
\section{The shape of the territory}\label{the-shape-of-the-territory}}

\begin{itemize}
\tightlist
\item
  \textbf{Composition} (ATL-01) supplies the arithmetic object.
\item
  The contrast between \textbf{finite and infinite alphabets} (ATL-02)
  fixes the \textbf{characteristic exponent} (ATL-03) and the
  \textbf{suppression} below it (ATL-04), which govern the
  \textbf{generated populations} (ATL-05).
\item
  At \textbf{fixed depth} the singularities sit at smaller abscissae
  (ATL-06): \textbf{depth two} is understood to its natural boundary
  (ATL-07), and the singularity at one third has a \textbf{local
  mechanism} whose general form is open (ATL-08).
\item
  \textbf{Collective growth} at and below \(s=1\) (ATL-09) concerns the
  passage to large depth where the sum over letters diverges.
\item
  \textbf{Extended generation} (ATL-10) adds multipliers that can
  deplete primes, and \textbf{entropy and admissibility} (ATL-11)
  describes the resulting exponent.
\end{itemize}

This list is a guide. The cited logical connections are given entry by
entry.

\hypertarget{atl-01-exact-admissible-composition}{%
\section{ATL-01: Exact admissible
composition}\label{atl-01-exact-admissible-composition}}

\hypertarget{mathematical-object}{%
\subsection{Mathematical object}\label{mathematical-object}}

A word of letters acting on integers, subject to integrality. The data
of a word are its minimal source \(K(w)\) and least endpoint \(E(w)\).

\hypertarget{established-structure}{%
\subsection{Established structure}\label{established-structure}}

\begin{itemize}
\tightlist
\item
  \textbf{Sources and endpoints are exact progressions.} A word is
  admissible precisely from the multiples of \(K(w)\), with endpoints
  the multiples of \(E(w)\), and both obey a one-letter gcd recurrence
  (T03). This holds for every integral-image grammar under the stated
  hypotheses.
\item
  \textbf{Prime coordinates.} Each letter shifts the exponent of \(p\)
  by \(2v_p(N+1)-v_p(N)\), and the minimal endpoint exponent is a
  maximum over suffixes (T04).
\item
  \textbf{Composition.} With \(d=\gcd(E(u),K(v))\) one has
  \(K(uv)=K(u)K(v)/d\) and \(E(uv)=E(u)E(v)/d\). Prime by prime this is
  multiplication in a bicyclic monoid, a classical object. It gives a
  block factorization of \(G_{m+n}\) and the supermultiplicativity
  \(G_{m+n}\ge G_mG_n\) (T47).
\item
  \textbf{History series.} Counting histories by endpoint factorizes as
  \(\mathcal Z_H(s,t)=\zeta(s)\,\mathcal G(s,t)\), with Möbius inversion
  depth by depth (T01, T02). A source weight gives a three-variable
  series with a closed form on its anti-diagonal (T14).
\item
  \textbf{No ordinary Euler product} exists for the nesting grammar at
  \(t=1\) (T06). Grammars with unit denominators, by contrast, have a
  closed rational form (T08).
\item
  At one step, the in-degree and the root series reduce to classical
  Dirichlet series (AS-001, AS-002).
\end{itemize}

\hypertarget{what-it-reveals}{%
\subsection{What it reveals}\label{what-it-reveals}}

Generation can be studied through two integers per word. The zeta factor
records only the free choice of source; the structure specific to the
grammar lies in \(\mathcal G\). Although the evolution of a prescribed
word separates into prime coordinates, the history population does not
possess an ordinary Euler product. The ordered letter structure couples
the arithmetic contributions when histories are counted.

\hypertarget{further-structure}{%
\subsection{Further structure}\label{further-structure}}

The recurrence of T03 drives the fixed-depth convergence region (ATL-06)
and the decompositions at one third (ATL-08).

\hypertarget{open-boundary}{%
\subsection{Open boundary}\label{open-boundary}}

The supermultiplicativity conjecture ANS-C2 for the history counts has
finite evidence only.

\hypertarget{reading-across-interpretation}{%
\subsection{Reading across
(interpretation)}\label{reading-across-interpretation}}

The prime-by-prime bicyclic structure is a point of contact with
algebraic investigations of the Substrate. It is examined in the
Cross-Structural Research Map (\texttt{CROSS\_STRUCTURES/}), whose
derived statements are unreviewed and are not part of this Atlas or the
catalogue.

\hypertarget{principal-results}{%
\subsection{Principal results}\label{principal-results}}

\begin{longtable}[]{@{}llll@{}}
\toprule
\begin{minipage}[b]{0.16\columnwidth}\raggedright
ID\strut
\end{minipage} & \begin{minipage}[b]{0.41\columnwidth}\raggedright
Result\strut
\end{minipage} & \begin{minipage}[b]{0.22\columnwidth}\raggedright
Status\strut
\end{minipage} & \begin{minipage}[b]{0.10\columnwidth}\raggedright
Review\strut
\end{minipage}\tabularnewline
\midrule
\endhead
\begin{minipage}[t]{0.16\columnwidth}\raggedright
ANS-T03\strut
\end{minipage} & \begin{minipage}[t]{0.41\columnwidth}\raggedright
Source and endpoint recurrences\strut
\end{minipage} & \begin{minipage}[t]{0.22\columnwidth}\raggedright
PROVED\strut
\end{minipage} & \begin{minipage}[t]{0.10\columnwidth}\raggedright
V\strut
\end{minipage}\tabularnewline
\begin{minipage}[t]{0.16\columnwidth}\raggedright
ANS-T04\strut
\end{minipage} & \begin{minipage}[t]{0.41\columnwidth}\raggedright
Prime-coordinate evolution\strut
\end{minipage} & \begin{minipage}[t]{0.22\columnwidth}\raggedright
PROVED\strut
\end{minipage} & \begin{minipage}[t]{0.10\columnwidth}\raggedright
V\strut
\end{minipage}\tabularnewline
\begin{minipage}[t]{0.16\columnwidth}\raggedright
ANS-T47\strut
\end{minipage} & \begin{minipage}[t]{0.41\columnwidth}\raggedright
Composition of words, block factorization, supermultiplicativity\strut
\end{minipage} & \begin{minipage}[t]{0.22\columnwidth}\raggedright
PROVED\strut
\end{minipage} & \begin{minipage}[t]{0.10\columnwidth}\raggedright
V\strut
\end{minipage}\tabularnewline
\begin{minipage}[t]{0.16\columnwidth}\raggedright
ANS-T01\strut
\end{minipage} & \begin{minipage}[t]{0.41\columnwidth}\raggedright
History Dirichlet factorization\strut
\end{minipage} & \begin{minipage}[t]{0.22\columnwidth}\raggedright
PROVED\strut
\end{minipage} & \begin{minipage}[t]{0.10\columnwidth}\raggedright
V\strut
\end{minipage}\tabularnewline
\begin{minipage}[t]{0.16\columnwidth}\raggedright
ANS-T06\strut
\end{minipage} & \begin{minipage}[t]{0.41\columnwidth}\raggedright
No ordinary Euler product\strut
\end{minipage} & \begin{minipage}[t]{0.22\columnwidth}\raggedright
PROVED\strut
\end{minipage} & \begin{minipage}[t]{0.10\columnwidth}\raggedright
V\strut
\end{minipage}\tabularnewline
\bottomrule
\end{longtable}

\hypertarget{recorded-open-questions}{%
\subsection{Recorded open questions}\label{recorded-open-questions}}

\begin{itemize}
\tightlist
\item
  \textbf{ANS-C2:} Conjecture (v001): H(mn) \textgreater{} H(m)H(n) for
  coprime non-squarefree m, n.~Finite evidence only.
\end{itemize}

\hypertarget{connected-entries-cited-dependencies}{%
\subsection{Connected entries (cited
dependencies)}\label{connected-entries-cited-dependencies}}

\begin{itemize}
\tightlist
\item
  \textbf{Uses ATL-03} (The characteristic exponent and its singular
  point): ANS-T07, ANS-T09.
\item
  \textbf{Uses ATL-06} (Fixed-depth singularity architecture): ANS-T12.
\item
  \textbf{Used by ATL-02} (Finite saturation and unbounded obstruction):
  ANS-T11.
\item
  \textbf{Used by ATL-06} (Fixed-depth singularity architecture):
  ANS-T12.
\item
  \textbf{Used by ATL-07} (Depth two: an arithmetic singularity
  hierarchy): ANS-T05, ANS-T13.
\item
  \textbf{Used by ATL-08} (The mechanism at one third): ANS-T46m,
  ANS-T46r, ANS-T46s.
\end{itemize}

\hypertarget{review-boundary}{%
\subsection{Review boundary}\label{review-boundary}}

Without a review verdict: AS-001 (B), AS-002 (B).

Qualified verdicts:

\begin{itemize}
\tightlist
\item
  ANS-T14: The stated convergence regions; not a complete classification
  for negative source weights.
\end{itemize}

\hypertarget{all-results-and-proof-sources}{%
\subsection{All results and proof
sources}\label{all-results-and-proof-sources}}

\begin{itemize}
\tightlist
\item
  \textbf{Initial analytic baseline:} AS-001, AS-002.
\item
  \textbf{ANS-I v001:} ANS-T03, ANS-T04, ANS-T01, ANS-T02, ANS-T14,
  ANS-T06, ANS-T08.
\item
  \textbf{ANS-I v014:} ANS-T47.
\end{itemize}

\hypertarget{atl-02-finite-saturation-and-unbounded-obstruction}{%
\section{ATL-02: Finite saturation and unbounded
obstruction}\label{atl-02-finite-saturation-and-unbounded-obstruction}}

\hypertarget{mathematical-object-1}{%
\subsection{Mathematical object}\label{mathematical-object-1}}

Random words with independent letters of law proportional to
\(r_N^{-a}\), with the alphabet either truncated at \(N\le M\) or
unrestricted, and the question whether an infinite word is admissible
from a finite integer.

\hypertarget{established-structure-1}{%
\subsection{Established structure}\label{established-structure-1}}

\begin{itemize}
\tightlist
\item
  \textbf{Every prime can be supplied from source 1}, and for every
  finite truncation \(1\le N\le M\) there is a finite preparation word
  after which the whole infinite random continuation is admissible with
  probability at least one half (T15).
\item
  \textbf{Criterion for other grammars.} Positive weighted drift at
  every consuming prime of a finite label set gives the corresponding
  lower growth bound (T11).
\item
  \textbf{At the finite critical exponent} \(\sigma_M\) the minimal
  source stabilizes almost surely, \(G_k^{(M)}(\sigma_M)\) decreases to
  a positive constant \(c_M\), and \(c_M\to0\) as \(M\to\infty\) (T26).
\item
  \textbf{With the full alphabet} and \(a>1\), a random infinite word is
  almost surely not admissible from any given integer (T17); hence
  \(K_k\to\infty\) almost surely (T18).
\end{itemize}

\hypertarget{what-it-reveals-1}{%
\subsection{What it reveals}\label{what-it-reveals-1}}

Every finite truncation can sustain indefinite random generation, while
unrestricted generation almost surely meets a demand it cannot supply.
The positive constants \(c_M\) vanishing in the limit measure how the
first regime passes into the second. The suppression theory of ATL-04
grows from this difference.

\hypertarget{further-structure-1}{%
\subsection{Further structure}\label{further-structure-1}}

The mechanisms differ. Saturation uses a positive drift in each of
finitely many prime coordinates. The obstruction uses fresh primes: a
prime not yet present is first met through a consuming letter with
probability bounded below, and this needs no independence between prime
coordinates.

\hypertarget{open-boundary-1}{%
\subsection{Open boundary}\label{open-boundary-1}}

A quantitative description of the transition from finite to infinite
alphabets (ANS-T30), and an analytic classification of grammars beyond
the saturation criterion (ANS-T11).

\hypertarget{reading-across-interpretation-1}{%
\subsection{Reading across
(interpretation)}\label{reading-across-interpretation-1}}

The pair suggests a test to apply to any alternative history grammar:
does positive drift in finitely many coordinates coexist with an
obstruction from unboundedly many?

\hypertarget{principal-results-1}{%
\subsection{Principal results}\label{principal-results-1}}

\begin{longtable}[]{@{}llll@{}}
\toprule
\begin{minipage}[b]{0.16\columnwidth}\raggedright
ID\strut
\end{minipage} & \begin{minipage}[b]{0.41\columnwidth}\raggedright
Result\strut
\end{minipage} & \begin{minipage}[b]{0.22\columnwidth}\raggedright
Status\strut
\end{minipage} & \begin{minipage}[b]{0.10\columnwidth}\raggedright
Review\strut
\end{minipage}\tabularnewline
\midrule
\endhead
\begin{minipage}[t]{0.16\columnwidth}\raggedright
ANS-T15\strut
\end{minipage} & \begin{minipage}[t]{0.41\columnwidth}\raggedright
Source-one generation: every prime reached; saturation\strut
\end{minipage} & \begin{minipage}[t]{0.22\columnwidth}\raggedright
PROVED\strut
\end{minipage} & \begin{minipage}[t]{0.10\columnwidth}\raggedright
V\strut
\end{minipage}\tabularnewline
\begin{minipage}[t]{0.16\columnwidth}\raggedright
ANS-T17\strut
\end{minipage} & \begin{minipage}[t]{0.41\columnwidth}\raggedright
Infinite-alphabet admissibility obstruction\strut
\end{minipage} & \begin{minipage}[t]{0.22\columnwidth}\raggedright
PROVED\strut
\end{minipage} & \begin{minipage}[t]{0.10\columnwidth}\raggedright
V\strut
\end{minipage}\tabularnewline
\begin{minipage}[t]{0.16\columnwidth}\raggedright
ANS-T18\strut
\end{minipage} & \begin{minipage}[t]{0.41\columnwidth}\raggedright
Divergence of minimal sources\strut
\end{minipage} & \begin{minipage}[t]{0.22\columnwidth}\raggedright
PROVED\strut
\end{minipage} & \begin{minipage}[t]{0.10\columnwidth}\raggedright
V\strut
\end{minipage}\tabularnewline
\begin{minipage}[t]{0.16\columnwidth}\raggedright
ANS-T26\strut
\end{minipage} & \begin{minipage}[t]{0.41\columnwidth}\raggedright
Finite-alphabet contrast\strut
\end{minipage} & \begin{minipage}[t]{0.22\columnwidth}\raggedright
PROVED\strut
\end{minipage} & \begin{minipage}[t]{0.10\columnwidth}\raggedright
V\strut
\end{minipage}\tabularnewline
\begin{minipage}[t]{0.16\columnwidth}\raggedright
ANS-T11\strut
\end{minipage} & \begin{minipage}[t]{0.41\columnwidth}\raggedright
Saturation criterion for general grammars (classification open)\strut
\end{minipage} & \begin{minipage}[t]{0.22\columnwidth}\raggedright
PROVED\strut
\end{minipage} & \begin{minipage}[t]{0.10\columnwidth}\raggedright
VQ\strut
\end{minipage}\tabularnewline
\bottomrule
\end{longtable}

\hypertarget{recorded-open-questions-1}{%
\subsection{Recorded open questions}\label{recorded-open-questions-1}}

\begin{itemize}
\tightlist
\item
  \textbf{ANS-T30:} Quantitative finite-to-infinite alphabet transition.
\item
  \textbf{ANS-T11:} Analytic classification across grammars (a
  saturation criterion is proved).
\end{itemize}

\hypertarget{connected-entries-cited-dependencies-1}{%
\subsection{Connected entries (cited
dependencies)}\label{connected-entries-cited-dependencies-1}}

\begin{itemize}
\tightlist
\item
  \textbf{Uses ATL-01} (Exact admissible composition): ANS-T03.
\item
  \textbf{Uses ATL-03} (The characteristic exponent and its singular
  point): ANS-T09.
\item
  \textbf{Uses ATL-04} (Collective critical suppression): ANS-T23.
\item
  \textbf{Used by ATL-03} (The characteristic exponent and its singular
  point): ANS-T09, ANS-T10, ANS-T19, ANS-T21, ANS-T31.
\item
  \textbf{Used by ATL-04} (Collective critical suppression): ANS-T23.
\end{itemize}

\hypertarget{review-boundary-1}{%
\subsection{Review boundary}\label{review-boundary-1}}

Every result in this entry has a review verdict.

Qualified verdicts:

\begin{itemize}
\tightlist
\item
  ANS-T11: The saturation criterion only; the general classification
  remains open.
\end{itemize}

\hypertarget{all-results-and-proof-sources-1}{%
\subsection{All results and proof
sources}\label{all-results-and-proof-sources-1}}

\begin{itemize}
\tightlist
\item
  \textbf{ANS-I v001:} ANS-T15, ANS-T11.
\item
  \textbf{ANS-I v002:} ANS-T26, ANS-T17, ANS-T18.
\end{itemize}

\hypertarget{atl-03-the-characteristic-exponent-and-its-singular-point}{%
\section{ATL-03: The characteristic exponent and its singular
point}\label{atl-03-the-characteristic-exponent-and-its-singular-point}}

\hypertarget{mathematical-object-2}{%
\subsection{Mathematical object}\label{mathematical-object-2}}

The full history series \(\mathcal G(s,1)=\sum_kG_k(s)\) and the
counting functions \(N_x(X)\), \(A(X)\), \(C(X)\).

\hypertarget{established-structure-2}{%
\subsection{Established structure}\label{established-structure-2}}

\begin{itemize}
\tightlist
\item
  \textbf{Growth in depth} is \(G_k(\sigma)^{1/k}\to B(\sigma)\) for
  \(\sigma>1\) and infinite for \(0\le\sigma\le1\), so the abscissa is
  \(\sigma_c\) (T09; a sufficient criterion is T07).
\item
  \textbf{Counting.} For every fixed source, and for the other two
  counts, \(\log Z(X)/\log X\to\sigma_c\) (T10).
\item
  \textbf{Probability form.} \(G_k(a)/B(a)^k=\mathbb E\,K_k^{-a}\)
  (T16). At \(\sigma_c\) the depth terms tend to zero although their
  \(k\)-th roots tend to one (T19).
\item
  \textbf{Below a pure power.} \(Z(X)=o(X^{\sigma_c})\) (T21).
\item
  \textbf{The singular point.} \(\mathcal G(s,1)\) has no meromorphic
  continuation through \(\sigma_c\) (T20), yet converges on the whole
  critical surface under the stated condition and is infinitely
  differentiable up to the line \(\operatorname{Re}s=\sigma_c\) (T22).
\item
  \textbf{The rest of the line is regular.} The function continues
  holomorphically across every point \(\sigma_c+i\tau\) with
  \(\tau\ne0\); the line is not a natural boundary (T31).
\end{itemize}

\hypertarget{what-it-reveals-2}{%
\subsection{What it reveals}\label{what-it-reveals-2}}

The exponent is the one that independent letters would give,
\(B(\sigma_c)=1\): admissibility does not move it. Admissibility changes
the nature of the singularity and pushes the counts below the pure
power. The singularity is smooth from the right and still not
meromorphic.

\hypertarget{further-structure-2}{%
\subsection{Further structure}\label{further-structure-2}}

The continuation of T31 comes from cancellation inside each \(G_k\); the
word series itself does not converge absolutely to the left of the line.

\hypertarget{open-boundary-2}{%
\subsection{Open boundary}\label{open-boundary-2}}

The nature of the singularity at \(\sigma_c\), including whether
continuations above and below it agree to its left (ANS-T33), and the
analogous theory for higher-degree and other nesting grammars (ANS-T25).

\hypertarget{reading-across-interpretation-2}{%
\subsection{Reading across
(interpretation)}\label{reading-across-interpretation-2}}

A smooth non-meromorphic singularity is a second kind of obstruction,
different in nature from the natural boundary of ATL-07.

\hypertarget{principal-results-2}{%
\subsection{Principal results}\label{principal-results-2}}

\begin{longtable}[]{@{}llll@{}}
\toprule
\begin{minipage}[b]{0.16\columnwidth}\raggedright
ID\strut
\end{minipage} & \begin{minipage}[b]{0.41\columnwidth}\raggedright
Result\strut
\end{minipage} & \begin{minipage}[b]{0.22\columnwidth}\raggedright
Status\strut
\end{minipage} & \begin{minipage}[b]{0.10\columnwidth}\raggedright
Review\strut
\end{minipage}\tabularnewline
\midrule
\endhead
\begin{minipage}[t]{0.16\columnwidth}\raggedright
ANS-T09\strut
\end{minipage} & \begin{minipage}[t]{0.41\columnwidth}\raggedright
Growth in depth and the sharp boundary sigma\_c\strut
\end{minipage} & \begin{minipage}[t]{0.22\columnwidth}\raggedright
PROVED\strut
\end{minipage} & \begin{minipage}[t]{0.10\columnwidth}\raggedright
V\strut
\end{minipage}\tabularnewline
\begin{minipage}[t]{0.16\columnwidth}\raggedright
ANS-T10\strut
\end{minipage} & \begin{minipage}[t]{0.41\columnwidth}\raggedright
Growth exponent of every counting function\strut
\end{minipage} & \begin{minipage}[t]{0.22\columnwidth}\raggedright
PROVED\strut
\end{minipage} & \begin{minipage}[t]{0.10\columnwidth}\raggedright
V\strut
\end{minipage}\tabularnewline
\begin{minipage}[t]{0.16\columnwidth}\raggedright
ANS-T20\strut
\end{minipage} & \begin{minipage}[t]{0.41\columnwidth}\raggedright
No meromorphic continuation through sigma\_c\strut
\end{minipage} & \begin{minipage}[t]{0.22\columnwidth}\raggedright
PROVED\strut
\end{minipage} & \begin{minipage}[t]{0.10\columnwidth}\raggedright
V\strut
\end{minipage}\tabularnewline
\begin{minipage}[t]{0.16\columnwidth}\raggedright
ANS-T22\strut
\end{minipage} & \begin{minipage}[t]{0.41\columnwidth}\raggedright
Convergence on the critical surface; smooth singularity\strut
\end{minipage} & \begin{minipage}[t]{0.22\columnwidth}\raggedright
PROVED\strut
\end{minipage} & \begin{minipage}[t]{0.10\columnwidth}\raggedright
V\strut
\end{minipage}\tabularnewline
\begin{minipage}[t]{0.16\columnwidth}\raggedright
ANS-T31\strut
\end{minipage} & \begin{minipage}[t]{0.41\columnwidth}\raggedright
Holomorphic continuation across the critical line away from
sigma\_c\strut
\end{minipage} & \begin{minipage}[t]{0.22\columnwidth}\raggedright
PROVED\strut
\end{minipage} & \begin{minipage}[t]{0.10\columnwidth}\raggedright
V\strut
\end{minipage}\tabularnewline
\begin{minipage}[t]{0.16\columnwidth}\raggedright
ANS-T21\strut
\end{minipage} & \begin{minipage}[t]{0.41\columnwidth}\raggedright
Strict growth suppression\strut
\end{minipage} & \begin{minipage}[t]{0.22\columnwidth}\raggedright
PROVED\strut
\end{minipage} & \begin{minipage}[t]{0.10\columnwidth}\raggedright
V\strut
\end{minipage}\tabularnewline
\bottomrule
\end{longtable}

\hypertarget{recorded-open-questions-2}{%
\subsection{Recorded open questions}\label{recorded-open-questions-2}}

\begin{itemize}
\tightlist
\item
  \textbf{ANS-T33:} Nature of the singularity at sigma\_c; whether the
  continuations above and below agree to its left.
\item
  \textbf{ANS-T25:} Higher-degree and other nesting grammars.
\end{itemize}

\hypertarget{connected-entries-cited-dependencies-2}{%
\subsection{Connected entries (cited
dependencies)}\label{connected-entries-cited-dependencies-2}}

\begin{itemize}
\tightlist
\item
  \textbf{Uses ATL-02} (Finite saturation and unbounded obstruction):
  ANS-T15, ANS-T18.
\item
  \textbf{Uses ATL-04} (Collective critical suppression): ANS-T23.
\item
  \textbf{Used by ATL-01} (Exact admissible composition): ANS-T14.
\item
  \textbf{Used by ATL-02} (Finite saturation and unbounded obstruction):
  ANS-T11.
\item
  \textbf{Used by ATL-04} (Collective critical suppression): ANS-T24.
\item
  \textbf{Used by ATL-05} (Generated populations and typical histories):
  ANS-T35, ANS-T36, ANS-T37, ANS-T37a, ANS-T40.
\end{itemize}

\hypertarget{review-boundary-2}{%
\subsection{Review boundary}\label{review-boundary-2}}

Every result in this entry has a review verdict.

\hypertarget{all-results-and-proof-sources-2}{%
\subsection{All results and proof
sources}\label{all-results-and-proof-sources-2}}

\begin{itemize}
\tightlist
\item
  \textbf{ANS-I v001:} ANS-T07, ANS-T09, ANS-T10.
\item
  \textbf{ANS-I v002:} ANS-T16, ANS-T19, ANS-T20, ANS-T21, ANS-T22.
\item
  \textbf{ANS-I v005:} ANS-T31.
\end{itemize}

\hypertarget{atl-04-collective-critical-suppression}{%
\section{ATL-04: Collective critical
suppression}\label{atl-04-collective-critical-suppression}}

\hypertarget{mathematical-object-3}{%
\subsection{Mathematical object}\label{mathematical-object-3}}

The negative moments \(\mathbb E\,K_k^{-b}\) of the minimal source, the
ratio \(Z(X)/X^{\sigma_c}\) and the derivatives of \(\mathcal G(s,1)\)
at \(\sigma_c\). Write \(\beta=1/a\), and \(\beta=1/\sigma_c\) at the
critical exponent.

\hypertarget{established-structure-3}{%
\subsection{Established structure}\label{established-structure-3}}

\begin{itemize}
\tightlist
\item
  \textbf{Depth.} For \(a>1\) and large \(k\),
  \(ck^{\beta}/\log k\le-\log\mathbb E\,K_k^{-b}\le Ck^{\beta}\) (T23a,
  sharpening T23).
\item
  \textbf{Counting.}
  \(c(\log X)^{\beta}/\log\log X\le-\log\big(Z(X)/X^{\sigma_c}\big)\le C(\log X)^{\beta}\)
  (T24a, sharpening T24; second proof of one inequality in T24b).
\item
  \textbf{Derivatives.} The \(j\)-th derivative magnitude \(D_j\)
  satisfies \(c^jj^{\sigma_cj}\le D_j\le(Cj\log j)^{\sigma_cj}\), so
  \(\sigma_c\) is the infimal Gevrey order on the closed half-plane
  (T27).
\item
  \textbf{First introductions.} The probability that every new prime
  among the first \(k\) letters arrives only through \(N+1\) has exact
  order \(\exp(-\Theta(k^{\beta}/\log k))\) (T28a).
\end{itemize}

\hypertarget{what-it-reveals-3}{%
\subsection{What it reveals}\label{what-it-reveals-3}}

The cost of admissibility at depth \(k\) is stretched-exponential, with
exponent fixed by the letter law. One quantity appears three times: as
decay in depth, as a deficit in counting and as the growth of
derivatives.

\hypertarget{further-structure-3}{%
\subsection{Further structure}\label{further-structure-3}}

The two sides of the depth estimate have different sources. That the
suppression is at least \(ck^{\beta}/\log k\) comes from fresh primes.
That it is at most \(Ck^{\beta}\) comes from Jensen's inequality and a
bound, for each prime, on the demand that can build up after the prime
is first consumed.

\hypertarget{open-boundary-3}{%
\subsection{Open boundary}\label{open-boundary-3}}

The exact order of the decay: the conjecture T28 states
\(-\log\mathbb E\,K_k^{-b}\asymp k^{\beta}/\log k\). T28a settles the
first introductions only. Preventing harmful first introductions does
not prevent depletion of primes met earlier, which is why T28a does not
resolve T28. Membership at exactly the Gevrey order \(\sigma_c\) is also
open (T27).

\hypertarget{reading-across-interpretation-3}{%
\subsection{Reading across
(interpretation)}\label{reading-across-interpretation-3}}

The gap between T28a and T28 isolates repeated demands on small primes.
The prime-resource questions of ATL-09 and ATL-11 have the same
character in other settings.

\hypertarget{principal-results-3}{%
\subsection{Principal results}\label{principal-results-3}}

\begin{longtable}[]{@{}llll@{}}
\toprule
\begin{minipage}[b]{0.16\columnwidth}\raggedright
ID\strut
\end{minipage} & \begin{minipage}[b]{0.41\columnwidth}\raggedright
Result\strut
\end{minipage} & \begin{minipage}[b]{0.22\columnwidth}\raggedright
Status\strut
\end{minipage} & \begin{minipage}[b]{0.10\columnwidth}\raggedright
Review\strut
\end{minipage}\tabularnewline
\midrule
\endhead
\begin{minipage}[t]{0.16\columnwidth}\raggedright
ANS-T23a\strut
\end{minipage} & \begin{minipage}[t]{0.41\columnwidth}\raggedright
Strengthened critical depth bounds\strut
\end{minipage} & \begin{minipage}[t]{0.22\columnwidth}\raggedright
PROVED\strut
\end{minipage} & \begin{minipage}[t]{0.10\columnwidth}\raggedright
V\strut
\end{minipage}\tabularnewline
\begin{minipage}[t]{0.16\columnwidth}\raggedright
ANS-T24a\strut
\end{minipage} & \begin{minipage}[t]{0.41\columnwidth}\raggedright
Growth suppression without logarithms in the lower estimate\strut
\end{minipage} & \begin{minipage}[t]{0.22\columnwidth}\raggedright
PROVED\strut
\end{minipage} & \begin{minipage}[t]{0.10\columnwidth}\raggedright
V\strut
\end{minipage}\tabularnewline
\begin{minipage}[t]{0.16\columnwidth}\raggedright
ANS-T27\strut
\end{minipage} & \begin{minipage}[t]{0.41\columnwidth}\raggedright
Critical derivative growth; infimal Gevrey order sigma\_c\strut
\end{minipage} & \begin{minipage}[t]{0.22\columnwidth}\raggedright
PROVED\strut
\end{minipage} & \begin{minipage}[t]{0.10\columnwidth}\raggedright
VQ\strut
\end{minipage}\tabularnewline
\begin{minipage}[t]{0.16\columnwidth}\raggedright
ANS-T28a\strut
\end{minipage} & \begin{minipage}[t]{0.41\columnwidth}\raggedright
Exact order of the first-introduction event\strut
\end{minipage} & \begin{minipage}[t]{0.22\columnwidth}\raggedright
PROVED\strut
\end{minipage} & \begin{minipage}[t]{0.10\columnwidth}\raggedright
V\strut
\end{minipage}\tabularnewline
\begin{minipage}[t]{0.16\columnwidth}\raggedright
ANS-T28\strut
\end{minipage} & \begin{minipage}[t]{0.41\columnwidth}\raggedright
Exact order of the critical decay\strut
\end{minipage} & \begin{minipage}[t]{0.22\columnwidth}\raggedright
CONJECTURE\strut
\end{minipage} & \begin{minipage}[t]{0.10\columnwidth}\raggedright
N\strut
\end{minipage}\tabularnewline
\bottomrule
\end{longtable}

\hypertarget{connected-entries-cited-dependencies-3}{%
\subsection{Connected entries (cited
dependencies)}\label{connected-entries-cited-dependencies-3}}

\begin{itemize}
\tightlist
\item
  \textbf{Uses ATL-02} (Finite saturation and unbounded obstruction):
  ANS-T15, ANS-T17.
\item
  \textbf{Uses ATL-03} (The characteristic exponent and its singular
  point): ANS-T21.
\item
  \textbf{Used by ATL-02} (Finite saturation and unbounded obstruction):
  ANS-T26.
\item
  \textbf{Used by ATL-03} (The characteristic exponent and its singular
  point): ANS-T22.
\item
  \textbf{Used by ATL-05} (Generated populations and typical histories):
  ANS-T37a, ANS-T37b, ANS-T42a.
\item
  \textbf{Used by ATL-09} (Collective growth at and below s = 1):
  ANS-T48.
\item
  \textbf{Used by ATL-10} (Extended generation and prime-resource
  balance): ANS2-T02, ANS2-T03, ANS2-T04, ANS2-T07, ANS2-T09a.
\item
  \textbf{Used by ATL-11} (Entropy and admissibility): ANS2N-T04b.
\end{itemize}

\hypertarget{review-boundary-3}{%
\subsection{Review boundary}\label{review-boundary-3}}

Without a review verdict: ANS-T28 (N).

Qualified verdicts:

\begin{itemize}
\tightlist
\item
  ANS-T27: Valid for the infimal Gevrey order; membership at exactly
  that order is open.
\end{itemize}

\hypertarget{all-results-and-proof-sources-3}{%
\subsection{All results and proof
sources}\label{all-results-and-proof-sources-3}}

\begin{itemize}
\tightlist
\item
  \textbf{ANS-I v002:} ANS-T23, ANS-T24.
\item
  \textbf{ANS-I v003:} ANS-T23a, ANS-T24a, ANS-T27, ANS-T28a, ANS-T28.
\item
  \textbf{ANS-I v004:} ANS-T24b.
\end{itemize}

\hypertarget{atl-05-generated-populations-and-typical-histories}{%
\section{ATL-05: Generated populations and typical
histories}\label{atl-05-generated-populations-and-typical-histories}}

\hypertarget{mathematical-object-4}{%
\subsection{Mathematical object}\label{mathematical-object-4}}

The histories counted up to endpoint \(X\), sorted by a homomorphism of
\(\mathbb Q_{>0}\) onto a finite abelian group \(\Gamma\) (the sector of
a word is the image of \(Q(w)\)).

\hypertarget{established-structure-4}{%
\subsection{Established structure}\label{established-structure-4}}

\begin{itemize}
\tightlist
\item
  \textbf{Analytic input.} Twisted series continue across the critical
  line except for the characters \(q\mapsto q^{i\tau}\) (T35); each
  sector function equals \(\mathcal G/|\Gamma|\) plus a function
  holomorphic across the line (T36).
\item
  \textbf{Equidistribution.} Critical log-moments equidistribute (T37a);
  \(A_\gamma(X)\sim A(X)/|\Gamma|\), and likewise from each fixed source
  (T37); at fixed depth the sector ratio at \(\sigma_c\) converges
  exponentially fast (T37b).
\item
  \textbf{Regular variation.} \(A(X)=X^{\sigma_c}L(X)\) with \(L\)
  slowly varying, and likewise for \(N_x\) (T40). The bounds of T24a
  apply to this factor.
\item
  \textbf{Replication.} \(C(X)\sim\zeta(\sigma_c)A(X)\) (T41).
\item
  \textbf{Typical history.} \(\log Q/\log X\to1\), \(\log K=o(\log X)\),
  depth over \(\log X\) tends to \(1/\lambda_c\), and block frequencies
  tend to the product weights \(r_N^{-\sigma_c}\) (T42), with
  concentration bounds (T42a).
\item
  \textbf{Limit laws.} The normalized endpoint has density
  \(\sigma_cu^{\sigma_c-1}\) (T43), and sector, replication factor and
  endpoint become independent (T44).
\end{itemize}

\hypertarget{what-it-reveals-4}{%
\subsection{What it reveals}\label{what-it-reveals-4}}

A typical generated history has the letter and block frequencies of an
independent sequence, although admissibility strongly suppresses the
total number of histories. The suppression acts on the size of the
population; the stated frequency laws are unaffected by it. This
establishes convergence of empirical block frequencies, rather than
independence of letters at arbitrary prescribed positions in the
conditioned history population.

\hypertarget{further-structure-4}{%
\subsection{Further structure}\label{further-structure-4}}

All of this passes through the holomorphic continuation of ATL-03 and
the characters of the multiplier group.

\hypertarget{open-boundary-4}{%
\subsection{Open boundary}\label{open-boundary-4}}

An asymptotic formula for the slowly varying factor is not available.
Twists \(q^{i\tau_0}\) at the matching height \(\tau=\tau_0\ne0\) are
open (ANS-T39).

\hypertarget{reading-across-interpretation-4}{%
\subsection{Reading across
(interpretation)}\label{reading-across-interpretation-4}}

The limit weights \(r_N^{-\sigma_c}\) describe a weighted history
population and are a candidate point of contact with generative
information geometry.

\hypertarget{principal-results-4}{%
\subsection{Principal results}\label{principal-results-4}}

\begin{longtable}[]{@{}llll@{}}
\toprule
ID & Result & Status & Review\tabularnewline
\midrule
\endhead
ANS-T37 & Counting equidistribution over finite quotients & PROVED &
V\tabularnewline
ANS-T40 & Regular variation with index sigma\_c & PROVED &
V\tabularnewline
ANS-T41 & Asymptotic arithmetic replication & PROVED & V\tabularnewline
ANS-T42 & Typical-history laws & PROVED & V\tabularnewline
ANS-T43 & Limit law of normalized endpoints & PROVED & V\tabularnewline
ANS-T44 & Joint limit law of sector, replication factor and endpoint &
PROVED & V\tabularnewline
\bottomrule
\end{longtable}

\hypertarget{recorded-open-questions-3}{%
\subsection{Recorded open questions}\label{recorded-open-questions-3}}

\begin{itemize}
\tightlist
\item
  \textbf{ANS-T39:} Twists q\^{}(i tau0) at tau = tau0 != 0.
\end{itemize}

\hypertarget{connected-entries-cited-dependencies-4}{%
\subsection{Connected entries (cited
dependencies)}\label{connected-entries-cited-dependencies-4}}

\begin{itemize}
\tightlist
\item
  \textbf{Uses ATL-03} (The characteristic exponent and its singular
  point): ANS-T10, ANS-T20, ANS-T22, ANS-T31.
\item
  \textbf{Uses ATL-04} (Collective critical suppression): ANS-T23a,
  ANS-T24a, ANS-T27.
\item
  \textbf{Uses ATL-06} (Fixed-depth singularity architecture): ANS-T12.
\end{itemize}

\hypertarget{review-boundary-4}{%
\subsection{Review boundary}\label{review-boundary-4}}

Every result in this entry has a review verdict.

\hypertarget{all-results-and-proof-sources-4}{%
\subsection{All results and proof
sources}\label{all-results-and-proof-sources-4}}

\begin{itemize}
\tightlist
\item
  \textbf{ANS-I v006:} ANS-T35, ANS-T36, ANS-T37a.
\item
  \textbf{ANS-I v007:} ANS-T37, ANS-T37b, ANS-T40, ANS-T41, ANS-T42.
\item
  \textbf{ANS-I v008:} ANS-T42a, ANS-T43, ANS-T44.
\end{itemize}

\hypertarget{atl-06-fixed-depth-singularity-architecture}{%
\section{ATL-06: Fixed-depth singularity
architecture}\label{atl-06-fixed-depth-singularity-architecture}}

\hypertarget{mathematical-object-5}{%
\subsection{Mathematical object}\label{mathematical-object-5}}

\(G_k(s)\) for a fixed depth \(k\), its restrictions to words whose
source is divisible by \(D\), and the counting function \(A_k(X)\).

\hypertarget{established-structure-5}{%
\subsection{Established structure}\label{established-structure-5}}

\begin{itemize}
\tightlist
\item
  \textbf{Region of convergence.} The weighted sum over words of length
  \(k\) converges exactly in the stated region; in particular \(G_k\)
  has abscissa \(1/2\) at every depth (T12).
\item
  \textbf{Counting.}
  \(A_k(X)\sim\kappa_k\sqrt X(\log X)^{k-1}/(2^{k-1}(k-1)!)\), with
  \(\kappa_k\) an Euler product (T32); as a corollary,
  \((2\sigma-1)^kG_k(\sigma)\to\kappa_k\) at \(1/2\) (C1). The local
  factors have the expansion
  \(1+\binom k2p^{-3/2}+\binom k3p^{-2}+O_k(p^{-5/2})\) (T32a).
\item
  \textbf{Continuation.} Every \(G_k\) is meromorphic on
  \(\operatorname{Re}s>1/3\), with a single pole there, at \(1/2\), of
  order exactly \(k\) (T34), also for every source restriction (T34b).
\item
  \textbf{Full main term.} \(A_k(X)=\sqrt X\,P_{k-1}(\log X)\) with a
  power saving (T34a).
\end{itemize}

\hypertarget{what-it-reveals-5}{%
\subsection{What it reveals}\label{what-it-reveals-5}}

At fixed depth \(k\), the principal pole has order \(k\), one order per
letter position. Arithmetic coupling changes its leading coefficient and
generates additional singularity structures.

\hypertarget{further-structure-5}{%
\subsection{Further structure}\label{further-structure-5}}

To the left of \(1/3\) the picture depends on depth: ATL-07 treats depth
two and ATL-08 the singularity at one third.

\hypertarget{open-boundary-5}{%
\subsection{Open boundary}\label{open-boundary-5}}

Zeros and poles of the complete functions (ANS-T38). Continuation of
general \(G_k\) to \(1/3\) and beyond is part of ATL-08.

\hypertarget{reading-across-interpretation-5}{%
\subsection{Reading across
(interpretation)}\label{reading-across-interpretation-5}}

The abscissae \(1/2\) at fixed depth and \(\sigma_c\) for the whole
series show that the infinite sum over depths creates a singularity that
no single depth has.

\hypertarget{principal-results-5}{%
\subsection{Principal results}\label{principal-results-5}}

\begin{longtable}[]{@{}llll@{}}
\toprule
ID & Result & Status & Review\tabularnewline
\midrule
\endhead
ANS-T12 & Convergence region at fixed depth & PROVED & V\tabularnewline
ANS-T32 & Fixed-depth counting asymptotic and the constants kappa\_k &
PROVED & VQ\tabularnewline
ANS-T34 & Every G\_k is meromorphic on Re s \textgreater{} 1/3 & PROVED
& V\tabularnewline
ANS-T34a & Full counting polynomial with power saving & PROVED &
V\tabularnewline
ANS-T34b & Exact pole order k at one half for every restriction & PROVED
& V\tabularnewline
\bottomrule
\end{longtable}

\hypertarget{recorded-open-questions-4}{%
\subsection{Recorded open questions}\label{recorded-open-questions-4}}

\begin{itemize}
\tightlist
\item
  \textbf{ANS-T38:} Complex zeros (and poles) of the complete functions
  G, G\_chi.
\end{itemize}

\hypertarget{connected-entries-cited-dependencies-5}{%
\subsection{Connected entries (cited
dependencies)}\label{connected-entries-cited-dependencies-5}}

\begin{itemize}
\tightlist
\item
  \textbf{Uses ATL-01} (Exact admissible composition): ANS-T03.
\item
  \textbf{Used by ATL-01} (Exact admissible composition): ANS-T47.
\item
  \textbf{Used by ATL-05} (Generated populations and typical histories):
  ANS-T35.
\item
  \textbf{Used by ATL-07} (Depth two: an arithmetic singularity
  hierarchy): ANS-T45.
\item
  \textbf{Used by ATL-08} (The mechanism at one third): ANS-T46k,
  ANS-T46m.
\end{itemize}

\hypertarget{review-boundary-5}{%
\subsection{Review boundary}\label{review-boundary-5}}

Every result in this entry has a review verdict.

Qualified verdicts:

\begin{itemize}
\tightlist
\item
  ANS-T32: Valid after a correction to the limiting step, which has been
  applied.
\end{itemize}

\hypertarget{all-results-and-proof-sources-5}{%
\subsection{All results and proof
sources}\label{all-results-and-proof-sources-5}}

\begin{itemize}
\tightlist
\item
  \textbf{ANS-I v001:} ANS-T12.
\item
  \textbf{ANS-I v005:} ANS-T32, ANS-C1.
\item
  \textbf{ANS-I v006:} ANS-T32a.
\item
  \textbf{ANS-I v008:} ANS-T34, ANS-T34a.
\item
  \textbf{ANS-I v009:} ANS-T34b.
\end{itemize}

\hypertarget{atl-07-depth-two-an-arithmetic-singularity-hierarchy}{%
\section{ATL-07: Depth two: an arithmetic singularity
hierarchy}\label{atl-07-depth-two-an-arithmetic-singularity-hierarchy}}

\hypertarget{mathematical-object-6}{%
\subsection{Mathematical object}\label{mathematical-object-6}}

\(G_2(s)=\sum_{a,b\ge1}\big[\gcd((a+1)^2,b)/((a+1)^2(b+1)^2)\big]^s\)
and the integer exponents \(e_n\) obtained by expanding its Euler factor
in powers of \(p^{-s}\).

\hypertarget{established-structure-6}{%
\subsection{Established structure}\label{established-structure-6}}

\begin{itemize}
\tightlist
\item
  \textbf{First continuation.} A Hurwitz-zeta representation gives
  continuation to \(\operatorname{Re}s>1/3\) and the double pole at
  \(1/2\), with \(A_2(X)\sim(\kappa/2)\sqrt X\log X\) (T05, T13).
\item
  \textbf{Factorization.} \(G_2\) is meromorphic on
  \(\operatorname{Re}s>1/5\) with a simple pole at \(1/3\) (T45), and on
  \(\operatorname{Re}s>0\), where the line \(\operatorname{Re}s=0\) is a
  natural boundary (T45a).
\item
  \textbf{Pole divisor.} On the positive real axis the poles are at
  \(1/n\) with \(e_n>0\), of order \(e_n\) (T45d). Examples: orders
  \(2,1,1,1\) at \(1/2,1/3,1/7,1/8\) (T45b) and order two at \(1/12\)
  (T45c).
\item
  \textbf{Exponents.} \(nc_n=-2R^n\cos n\theta+O(\tau(n)R^{n/2})\) for
  the stated algebraic \(R>1\) and \(\theta\), each sign of \(e_n\) has
  density one half (T45e), and the normalized exponents follow an
  arcsine law (T45f).
\item
  \textbf{Zeta zeros.} For \(e_n<0\), all but a small set of
  critical-line zeros \(\rho\) give poles at \(\rho/n\) (T38d; cases
  T38a, T38b), with lower bounds (T38c) and asymptotic populations on
  the lines \(\operatorname{Re}s=1/(2n)\) (T38e).
\end{itemize}

The description of \textbf{all} non-real poles in the second corollary
of T45d assumes the Riemann hypothesis. The other statements listed here
are recorded without it.

\hypertarget{what-it-reveals-6}{%
\subsection{What it reveals}\label{what-it-reveals-6}}

A single coupling between two letters produces an unbounded hierarchy of
singularities, accumulating on a natural boundary, with orders governed
by one cubic polynomial.

\hypertarget{further-structure-6}{%
\subsection{Further structure}\label{further-structure-6}}

The signs of the exponents decide whether a given \(n\) contributes
poles on the real axis or poles at rescaled zeta zeros.

\hypertarget{open-boundary-6}{%
\subsection{Open boundary}\label{open-boundary-6}}

The maximal meromorphic domain and the pole divisor are established. Two
questions about depth two remain unresolved: the classification of the
zeros of \(G_2\) (ANS-G2Z), and a rigorous sharp, unsmoothed secondary
counting term of order \(X^{1/3}\) in \(A_2(X)\) (ANS-A2S). They are
preserved as mathematical questions, not active investigations. The
corresponding hierarchy at higher depth is also open (ATL-08).

\hypertarget{reading-across-interpretation-6}{%
\subsection{Reading across
(interpretation)}\label{reading-across-interpretation-6}}

The natural boundary is a candidate point of contact with structural
obstruction theory. The poles at rescaled zeta zeros reflect zeta
factors in the formula; they carry no information about where those
zeros lie.

\hypertarget{principal-results-6}{%
\subsection{Principal results}\label{principal-results-6}}

\begin{longtable}[]{@{}llll@{}}
\toprule
\begin{minipage}[b]{0.16\columnwidth}\raggedright
ID\strut
\end{minipage} & \begin{minipage}[b]{0.41\columnwidth}\raggedright
Result\strut
\end{minipage} & \begin{minipage}[b]{0.22\columnwidth}\raggedright
Status\strut
\end{minipage} & \begin{minipage}[b]{0.10\columnwidth}\raggedright
Review\strut
\end{minipage}\tabularnewline
\midrule
\endhead
\begin{minipage}[t]{0.16\columnwidth}\raggedright
ANS-T45\strut
\end{minipage} & \begin{minipage}[t]{0.41\columnwidth}\raggedright
Depth two to Re s \textgreater{} 1/5 with a simple pole at 1/3\strut
\end{minipage} & \begin{minipage}[t]{0.22\columnwidth}\raggedright
PROVED\strut
\end{minipage} & \begin{minipage}[t]{0.10\columnwidth}\raggedright
V\strut
\end{minipage}\tabularnewline
\begin{minipage}[t]{0.16\columnwidth}\raggedright
ANS-T45a\strut
\end{minipage} & \begin{minipage}[t]{0.41\columnwidth}\raggedright
Depth two: meromorphic on Re s \textgreater{} 0 with natural
boundary\strut
\end{minipage} & \begin{minipage}[t]{0.22\columnwidth}\raggedright
PROVED\strut
\end{minipage} & \begin{minipage}[t]{0.10\columnwidth}\raggedright
V\strut
\end{minipage}\tabularnewline
\begin{minipage}[t]{0.16\columnwidth}\raggedright
ANS-T45d\strut
\end{minipage} & \begin{minipage}[t]{0.41\columnwidth}\raggedright
Depth two: exact pole divisor on Re s \textgreater{} 0\strut
\end{minipage} & \begin{minipage}[t]{0.22\columnwidth}\raggedright
PROVED\strut
\end{minipage} & \begin{minipage}[t]{0.10\columnwidth}\raggedright
V\strut
\end{minipage}\tabularnewline
\begin{minipage}[t]{0.16\columnwidth}\raggedright
ANS-T45e\strut
\end{minipage} & \begin{minipage}[t]{0.41\columnwidth}\raggedright
Depth two: asymptotics and sign densities of the Euler exponents\strut
\end{minipage} & \begin{minipage}[t]{0.22\columnwidth}\raggedright
PROVED\strut
\end{minipage} & \begin{minipage}[t]{0.10\columnwidth}\raggedright
V\strut
\end{minipage}\tabularnewline
\begin{minipage}[t]{0.16\columnwidth}\raggedright
ANS-T38d\strut
\end{minipage} & \begin{minipage}[t]{0.41\columnwidth}\raggedright
Depth two: vertical pole family for every negative exponent\strut
\end{minipage} & \begin{minipage}[t]{0.22\columnwidth}\raggedright
PROVED\strut
\end{minipage} & \begin{minipage}[t]{0.10\columnwidth}\raggedright
V\strut
\end{minipage}\tabularnewline
\begin{minipage}[t]{0.16\columnwidth}\raggedright
ANS-T38e\strut
\end{minipage} & \begin{minipage}[t]{0.41\columnwidth}\raggedright
Depth two: pole populations on the lines Re s = 1/(2n)\strut
\end{minipage} & \begin{minipage}[t]{0.22\columnwidth}\raggedright
PROVED\strut
\end{minipage} & \begin{minipage}[t]{0.10\columnwidth}\raggedright
V\strut
\end{minipage}\tabularnewline
\bottomrule
\end{longtable}

\hypertarget{recorded-open-questions-5}{%
\subsection{Recorded open questions}\label{recorded-open-questions-5}}

\begin{itemize}
\tightlist
\item
  \textbf{ANS-G2Z:} Classification of the zeros of G\_2 (the maximal
  meromorphic domain and the pole divisor are established). ID is a
  label of this collection.
\item
  \textbf{ANS-A2S:} Rigorous sharp, unsmoothed secondary counting term
  of order X\^{}(1/3) in A\_2(X) (recorded as not proved in ANS-I v008;
  a smoothed form holds). ID is a label of this collection.
\end{itemize}

\hypertarget{connected-entries-cited-dependencies-6}{%
\subsection{Connected entries (cited
dependencies)}\label{connected-entries-cited-dependencies-6}}

\begin{itemize}
\tightlist
\item
  \textbf{Uses ATL-01} (Exact admissible composition): ANS-T03.
\item
  \textbf{Uses ATL-06} (Fixed-depth singularity architecture): ANS-T34.
\item
  \textbf{Used by ATL-08} (The mechanism at one third): ANS-T46b,
  ANS-T46h-i, ANS-T46k, ANS-T46s.
\end{itemize}

\hypertarget{review-boundary-6}{%
\subsection{Review boundary}\label{review-boundary-6}}

Every result in this entry has a review verdict.

\hypertarget{all-results-and-proof-sources-6}{%
\subsection{All results and proof
sources}\label{all-results-and-proof-sources-6}}

\begin{itemize}
\tightlist
\item
  \textbf{ANS-I v001:} ANS-T13, ANS-T05.
\item
  \textbf{ANS-I v008:} ANS-T45.
\item
  \textbf{ANS-I v009:} ANS-T45a, ANS-T45b, ANS-T38a.
\item
  \textbf{ANS-I v010:} ANS-T45d, ANS-T45c, ANS-T45e, ANS-T38d, ANS-T38b,
  ANS-T38c.
\item
  \textbf{ANS-I v011:} ANS-T45f, ANS-T38e.
\end{itemize}

\hypertarget{atl-08-the-mechanism-at-one-third}{%
\section{ATL-08: The mechanism at one
third}\label{atl-08-the-mechanism-at-one-third}}

\hypertarget{mathematical-object-7}{%
\subsection{Mathematical object}\label{mathematical-object-7}}

The singularity of \(G_k\) at \(s=1/3\), produced by primes shared
between an earlier letter through \(N+1\) and a later letter through
\(N\).

\hypertarget{established-structure-depth-three}{%
\subsection{Established structure: depth
three}\label{established-structure-depth-three}}

\begin{itemize}
\tightlist
\item
  An exact decomposition expresses \(G_3\) through a shifted series
  \(V\), a remainder \(R\) and \(G_2\) (T46d). Models of one linked
  letter have only a simple pole at \(1/3\) (T46c, T46f).
\item
  Square-divisor identities (T46g) lead to the continuation of the
  remainder (T46h, T46i), whose first component has a natural boundary
  at \(\operatorname{Re}s=0\) (T46h-i). A theorem on shifted series
  (T46j) then gives the continuation and residue of \(V\) (T46e);
  positivity of its constants was proved in advance (T46e-i).
\item
  \(G_3\) has a pole of order \textbf{exactly two} at \(1/3\), with an
  explicit leading coefficient near \(-0.75916\) (T46k); the
  continuation reaches \(\operatorname{Re}s>1/3-1/200\) (T46k-i), and
  smoothed counts have a second term of size \(X^{1/3}\log X\) (T46l).
\end{itemize}

\hypertarget{established-structure-arbitrary-depth}{%
\subsection{Established structure: arbitrary
depth}\label{established-structure-arbitrary-depth}}

\begin{itemize}
\tightlist
\item
  For every \(k\ge2\), on \(\operatorname{Re}s>1/2\), \(G_k\) decomposes
  by its end letters into a weighted middle-word sum \(M_k\), a
  remainder and \(G_{k-1}\) (T46m). Overlaps in the weight obey an exact
  identity (T46r).
\item
  At one prime at most \(\lfloor k^2/4\rfloor\) links can be active
  (T46n). A model with \(r\) supplying and \(t\) receiving letters has a
  pole of order exactly \(rt\) (T46o).
\item
  For a role pattern \(R\) the local expectation is
  \(1+L(R)/p+O(p^{-4/3})\), where \(L(R)\) counts links; the resulting
  formula reproduces the proved leading coefficients at \(k=2,3\) and in
  the model (T46s).
\item
  Stated conditions on \(M_k\), the remainder and \(G_{k-1}\) would
  imply the pole order at depth \(k\) (T46p). They are not established
  for any \(k\ge4\).
\end{itemize}

\hypertarget{what-it-reveals-7}{%
\subsection{What it reveals}\label{what-it-reveals-7}}

Under the role-pattern measure, each eligible supplier-to-receiver link
contributes to the first-order coefficient \(L(R)/p\) of the local
expectation. The maximal link count is attained by a balanced complete
bipartite pattern, with all suppliers preceding all receivers. This
mechanism reproduces the proved singularities at depths two and three
and in the comparison model. Its extension to the full history series at
arbitrary depth remains conjectural.

\hypertarget{further-structure-7}{%
\subsection{Further structure}\label{further-structure-7}}

Three statements of different strength are kept apart:

\begin{longtable}[]{@{}lll@{}}
\toprule
\begin{minipage}[b]{0.18\columnwidth}\raggedright
Identifier\strut
\end{minipage} & \begin{minipage}[b]{0.58\columnwidth}\raggedright
Content\strut
\end{minipage} & \begin{minipage}[b]{0.16\columnwidth}\raggedright
Proved\strut
\end{minipage}\tabularnewline
\midrule
\endhead
\begin{minipage}[t]{0.18\columnwidth}\raggedright
T46b\strut
\end{minipage} & \begin{minipage}[t]{0.58\columnwidth}\raggedright
\(G_k\) has a pole at \(1/3\) of order \(\lfloor k^2/4\rfloor\)\strut
\end{minipage} & \begin{minipage}[t]{0.16\columnwidth}\raggedright
\(k=2,3\)\strut
\end{minipage}\tabularnewline
\begin{minipage}[t]{0.18\columnwidth}\raggedright
T46q\strut
\end{minipage} & \begin{minipage}[t]{0.58\columnwidth}\raggedright
\(M_k\) has a pole of order exactly \(\lfloor k^2/4\rfloor-1\)\strut
\end{minipage} & \begin{minipage}[t]{0.16\columnwidth}\raggedright
\(k=2,3\)\strut
\end{minipage}\tabularnewline
\begin{minipage}[t]{0.18\columnwidth}\raggedright
T46t\strut
\end{minipage} & \begin{minipage}[t]{0.58\columnwidth}\raggedright
Continuation near \(1/3\), pole order \(\lfloor k^2/4\rfloor\) and the
leading coefficient\strut
\end{minipage} & \begin{minipage}[t]{0.16\columnwidth}\raggedright
\(k=2,3\)\strut
\end{minipage}\tabularnewline
\bottomrule
\end{longtable}

In the second row, order zero at \(k=2\) denotes regularity, not a pole:
\(M_2=1\) is regular and nonzero at \(\frac13\).

\hypertarget{open-boundary-7}{%
\subsection{Open boundary}\label{open-boundary-7}}

Depth four and beyond at and past \(\operatorname{Re}s=1/3\), and depth
three further left (ANS-T46); singularities and natural boundaries at
higher depth (ANS-T46a). ANS-I v014 records that no viable route to the
continuation of \(G_4\) past \(\frac13\) has been established.

\hypertarget{reading-across-interpretation-7}{%
\subsection{Reading across
(interpretation)}\label{reading-across-interpretation-7}}

The quadratic growth of the conjectured order with depth is a candidate
point of contact with generative geometry. Models are not the functions
\(G_k\).

\hypertarget{principal-results-7}{%
\subsection{Principal results}\label{principal-results-7}}

\begin{longtable}[]{@{}llll@{}}
\toprule
\begin{minipage}[b]{0.16\columnwidth}\raggedright
ID\strut
\end{minipage} & \begin{minipage}[b]{0.41\columnwidth}\raggedright
Result\strut
\end{minipage} & \begin{minipage}[b]{0.22\columnwidth}\raggedright
Status\strut
\end{minipage} & \begin{minipage}[b]{0.10\columnwidth}\raggedright
Review\strut
\end{minipage}\tabularnewline
\midrule
\endhead
\begin{minipage}[t]{0.16\columnwidth}\raggedright
ANS-T46k\strut
\end{minipage} & \begin{minipage}[t]{0.41\columnwidth}\raggedright
Depth three: pole of order exactly two at 1/3\strut
\end{minipage} & \begin{minipage}[t]{0.22\columnwidth}\raggedright
PROVED\strut
\end{minipage} & \begin{minipage}[t]{0.10\columnwidth}\raggedright
V\strut
\end{minipage}\tabularnewline
\begin{minipage}[t]{0.16\columnwidth}\raggedright
ANS-T46m\strut
\end{minipage} & \begin{minipage}[t]{0.41\columnwidth}\raggedright
Every depth: decomposition by end letters\strut
\end{minipage} & \begin{minipage}[t]{0.22\columnwidth}\raggedright
PROVED\strut
\end{minipage} & \begin{minipage}[t]{0.10\columnwidth}\raggedright
V\strut
\end{minipage}\tabularnewline
\begin{minipage}[t]{0.16\columnwidth}\raggedright
ANS-T46s\strut
\end{minipage} & \begin{minipage}[t]{0.41\columnwidth}\raggedright
Role patterns: local expectation and agreement with proved
coefficients\strut
\end{minipage} & \begin{minipage}[t]{0.22\columnwidth}\raggedright
PROVED\strut
\end{minipage} & \begin{minipage}[t]{0.10\columnwidth}\raggedright
V\strut
\end{minipage}\tabularnewline
\begin{minipage}[t]{0.16\columnwidth}\raggedright
ANS-T46b\strut
\end{minipage} & \begin{minipage}[t]{0.41\columnwidth}\raggedright
Pole of order floor(k\^{}2/4) at 1/3\strut
\end{minipage} & \begin{minipage}[t]{0.22\columnwidth}\raggedright
PROVED for k \textless= 3; CONJECTURE for k \textgreater= 4\strut
\end{minipage} & \begin{minipage}[t]{0.10\columnwidth}\raggedright
V\strut
\end{minipage}\tabularnewline
\begin{minipage}[t]{0.16\columnwidth}\raggedright
ANS-T46q\strut
\end{minipage} & \begin{minipage}[t]{0.41\columnwidth}\raggedright
Weighted middle-word sum: pole order floor(k\^{}2/4)-1\strut
\end{minipage} & \begin{minipage}[t]{0.22\columnwidth}\raggedright
CONJECTURE (k = 2, 3 proved)\strut
\end{minipage} & \begin{minipage}[t]{0.10\columnwidth}\raggedright
N\strut
\end{minipage}\tabularnewline
\begin{minipage}[t]{0.16\columnwidth}\raggedright
ANS-T46t\strut
\end{minipage} & \begin{minipage}[t]{0.41\columnwidth}\raggedright
General continuation, pole order floor(k\^{}2/4) and leading coefficient
at 1/3\strut
\end{minipage} & \begin{minipage}[t]{0.22\columnwidth}\raggedright
CONJECTURE (k = 2, 3 proved)\strut
\end{minipage} & \begin{minipage}[t]{0.10\columnwidth}\raggedright
N\strut
\end{minipage}\tabularnewline
\bottomrule
\end{longtable}

\hypertarget{recorded-open-questions-6}{%
\subsection{Recorded open questions}\label{recorded-open-questions-6}}

\begin{itemize}
\tightlist
\item
  \textbf{ANS-T46:} Depth k \textgreater= 4 at and beyond Re s = 1/3;
  depth three further to the left.
\item
  \textbf{ANS-T46a:} Singularities and natural boundaries at higher
  depth.
\end{itemize}

\hypertarget{connected-entries-cited-dependencies-7}{%
\subsection{Connected entries (cited
dependencies)}\label{connected-entries-cited-dependencies-7}}

\begin{itemize}
\tightlist
\item
  \textbf{Uses ATL-01} (Exact admissible composition): ANS-T03, ANS-T47.
\item
  \textbf{Uses ATL-06} (Fixed-depth singularity architecture): ANS-T12,
  ANS-T34.
\item
  \textbf{Uses ATL-07} (Depth two: an arithmetic singularity hierarchy):
  ANS-T45, ANS-T45a.
\end{itemize}

\hypertarget{review-boundary-7}{%
\subsection{Review boundary}\label{review-boundary-7}}

Without a review verdict: ANS-T46q (N), ANS-T46t (N).

Qualified verdicts:

\begin{itemize}
\tightlist
\item
  ANS-T46h-i: A normal-convergence clarification was requested and added
  in ANS-I v013.
\item
  ANS-T46j: Two scope clarifications were requested and made in ANS-I
  v013.
\end{itemize}

\hypertarget{all-results-and-proof-sources-7}{%
\subsection{All results and proof
sources}\label{all-results-and-proof-sources-7}}

\begin{itemize}
\tightlist
\item
  \textbf{ANS-I v010:} ANS-T46c.
\item
  \textbf{ANS-I v011:} ANS-T46f, ANS-T46d, ANS-T46e-i.
\item
  \textbf{ANS-I v012:} ANS-T46g, ANS-T46h, ANS-T46h-i, ANS-T46i,
  ANS-T46j, ANS-T46e, ANS-T46k, ANS-T46l.
\item
  \textbf{ANS-I v013:} ANS-T46k-i, ANS-T46m, ANS-T46n, ANS-T46o,
  ANS-T46p, ANS-T46q.
\item
  \textbf{ANS-I v014:} ANS-T46r, ANS-T46s.
\item
  \textbf{ANS-I v009:} ANS-T46b.
\item
  \textbf{ANS-I v015:} ANS-T46t.
\end{itemize}

\hypertarget{atl-09-collective-growth-at-and-below-s-1}{%
\section{ATL-09: Collective growth at and below s =
1}\label{atl-09-collective-growth-at-and-below-s-1}}

\hypertarget{mathematical-object-8}{%
\subsection{Mathematical object}\label{mathematical-object-8}}

\(G_k(s)^{1/k}\) for real \(s>1/2\), with particular attention to
\(1/2<s\le1\) and the window near \(1\). For \(s\le1\) one has
\(B(s)=\infty\) and the radius of convergence of \(\sum_kG_k(s)t^k\) is
zero. For \(s\le1/2\) the fixed-depth sums already diverge.

\hypertarget{established-structure-7}{%
\subsection{Established structure}\label{established-structure-7}}

\begin{itemize}
\tightlist
\item
  At \(s=1\) the growth has exact order \(\log k\):
  \(\log k-\log\log k+O(1)\le G_k(1)^{1/k}\le(e^{\gamma}\zeta(2)+o(1))\log k\),
  and for \(1/2<s<1\) it is at least \(c_sk^{1-s}\) (T48).
\item
  The constants are sharpened on both sides (T49 below, T48a above):
  \[\boxed{1.107910\ldots\le\liminf_{k\to\infty}\frac{G_k(1)^{1/k}}{\log k}\le\limsup_{k\to\infty}\frac{G_k(1)^{1/k}}{\log k}\le e^{\gamma}}\]
\item
  Under growth constants \(0<c_-\le c_+<\infty\), the exponential
  generating function \(\sum_kG_k(1)t^k/k!\) is entire of order one and
  infinite type, with matching maximum-modulus bounds (T48b).
\item
  Two-sided bounds hold in the window \(s=1+\tau/\log k\) for every
  fixed real \(\tau\) as \(k\to\infty\) (T50). No uniform statement for
  a growing parameter \(\tau=\tau(k)\) is established.
\end{itemize}

\hypertarget{what-it-reveals-8}{%
\subsection{What it reveals}\label{what-it-reveals-8}}

Without admissibility the value at \(s=1\) would be infinite. With it,
the growth rate is finite and of size \(\log k\), which is also the
total weight of the letters up to a fixed power of \(k\).

\hypertarget{further-structure-8}{%
\subsection{Further structure}\label{further-structure-8}}

The lower constant above one uses letters beyond the base range whose
prime factors are small, counted with the Dickman function (T49).

\hypertarget{open-boundary-8}{%
\subsection{Open boundary}\label{open-boundary-8}}

Existence and value of \(\lim G_k(1)^{1/k}/\log k\) (ANS-LIM), including
whether the upper limit is strictly below \(e^{\gamma}\). A heuristic
recorded in v015, which assumes an equidistribution that is not proved,
suggests a value below \(e^{\gamma}\).

\hypertarget{reading-across-interpretation-8}{%
\subsection{Reading across
(interpretation)}\label{reading-across-interpretation-8}}

The constraint here is a balance between primes supplied by \(N+1\) and
consumed by \(N\), the same bookkeeping as in ATL-10. The variational
principle of ATL-11 gives a broader context for this reading. No theorem
in the collection carries it over to \(s=1\): those results assume
\(a>1\) in the extended grammar.

\hypertarget{principal-results-8}{%
\subsection{Principal results}\label{principal-results-8}}

\begin{longtable}[]{@{}llll@{}}
\toprule
ID & Result & Status & Review\tabularnewline
\midrule
\endhead
ANS-T48 & Growth of G\_k at s = 1 of exact order log k & PROVED &
V\tabularnewline
ANS-T48a & Upper constant e\^{}gamma & PROVED & V\tabularnewline
ANS-T49 & Lower constant 1.1079 & PROVED & V\tabularnewline
ANS-T50 & Growth bounds in the window around s = 1 & PROVED &
VQ\tabularnewline
\bottomrule
\end{longtable}

\hypertarget{recorded-open-questions-7}{%
\subsection{Recorded open questions}\label{recorded-open-questions-7}}

\begin{itemize}
\tightlist
\item
  \textbf{ANS-LIM:} Existence and value of lim G\_k(1)\^{}(1/k)/log k. A
  heuristic (v015, unproved) suggests a value below e\^{}gamma. No
  separate ID in the releases; ANS-LIM is a label of this collection.
\end{itemize}

\hypertarget{connected-entries-cited-dependencies-8}{%
\subsection{Connected entries (cited
dependencies)}\label{connected-entries-cited-dependencies-8}}

\begin{itemize}
\tightlist
\item
  \textbf{Uses ATL-04} (Collective critical suppression): ANS-T23,
  ANS-T23a.
\end{itemize}

\hypertarget{review-boundary-8}{%
\subsection{Review boundary}\label{review-boundary-8}}

Every result in this entry has a review verdict.

Qualified verdicts:

\begin{itemize}
\tightlist
\item
  ANS-T48b: The bounding constants satisfy 0 \textless{} c\_- \textless=
  c\_+ \textless{} infinity.
\item
  ANS-T50: Fixed real tau as k tends to infinity; no uniform statement
  for tau = tau(k).
\end{itemize}

\hypertarget{all-results-and-proof-sources-8}{%
\subsection{All results and proof
sources}\label{all-results-and-proof-sources-8}}

\begin{itemize}
\tightlist
\item
  \textbf{ANS-I v014:} ANS-T48.
\item
  \textbf{ANS-I v015:} ANS-T48a, ANS-T48b, ANS-T49, ANS-T50.
\end{itemize}

\hypertarget{atl-10-extended-generation-and-prime-resource-balance}{%
\section{ATL-10: Extended generation and prime-resource
balance}\label{atl-10-extended-generation-and-prime-resource-balance}}

\hypertarget{mathematical-object-9}{%
\subsection{Mathematical object}\label{mathematical-object-9}}

The nesting grammar with finitely many further multipliers
\(\alpha_j>1\) having a rational power. States fall into finitely many
torsion classes. The \textbf{drift} \(\mu_p(a)\) is the mean change of
the exponent of \(p\) per letter under the law proportional to
\(\mu(\ell)^{-a}\); \(\sigma_A\) solves \(B_A(\sigma_A)=1\).

\hypertarget{established-structure-8}{%
\subsection{Established structure}\label{established-structure-8}}

\begin{itemize}
\tightlist
\item
  \textbf{Foundations.} Source formula, comparison of classes,
  factorization \(\mathcal Z_\varphi=\zeta\,\mathcal G_\varphi\) and
  invariance on valuation-type classes (L1, L2, T01, T00).
\item
  \textbf{Positive drift.} Depth bounds, counting bounds, derivative
  bounds and infimal Gevrey order \(\sigma_A\) transfer from the nesting
  theory (T02, T03, T04); first introductions keep their order (T07).
  The exponent depends on the real values of the multipliers (T08).
\item
  \textbf{Operators.} A first-letter identity (T06), one-letter formulae
  for any signs and a closed countable operator system for
  \(\operatorname{Re}s>1\), with its realizations (T14, T10 to T13,
  T15).
\item
  \textbf{Growth rate in general.}
  \(\lim\mathcal G_{\varphi,k}(a)^{1/k}=\rho_A(a)\), the minimum of an
  explicit convex function. The abscissa \(\sigma^\ast\) solves
  \(\rho_A=1\), is the same in every class, equals \(\sigma_A\) if all
  drifts at \(\sigma_A\) are nonnegative and is strictly smaller if one
  is negative (T16; earlier bound T05).
\item
  \textbf{Boundary cases.} If some drifts vanish at \(\sigma_A\) and the
  others are positive, the series is finite at the abscissa (T09). If
  all drifts are nonnegative, depth bounds of the form
  \(ck^{\beta}/\log k\) to \(Ck^{\beta}\) still hold (T09a).
\item
  \textbf{Example.} An interval for the abscissa with multiplier
  \(2/\sqrt3\) (C16, floating point; certified as ANS2N-T06).
\end{itemize}

\hypertarget{what-it-reveals-9}{%
\subsection{What it reveals}\label{what-it-reveals-9}}

The sign of the drift decides whether admissibility changes only
lower-order behaviour or changes the exponent itself. If no drift is
negative, the exponent stays at the characteristic root \(\sigma_A\), as
in the original grammar; if some prime is depleted on average, it lies
strictly below that root.

\hypertarget{further-structure-9}{%
\subsection{Further structure}\label{further-structure-9}}

T16 expresses the effect of depletion as a minimization over prime
weights, developed in ATL-11.

\hypertarget{open-boundary-9}{%
\subsection{Open boundary}\label{open-boundary-9}}

Exact suppression order in the extensions, behaviour at the abscissa
under negative drift, the spectrum and continuation of the transfer
operators, and whether a drift can vanish exactly at \(\sigma_A\)
(ANS2-O3, O5, O6, O7).

\hypertarget{reading-across-interpretation-9}{%
\subsection{Reading across
(interpretation)}\label{reading-across-interpretation-9}}

The dichotomy is a candidate point of contact with admissibility and
arithmetic dynamics, and with the study of alternative history grammars.

\hypertarget{principal-results-9}{%
\subsection{Principal results}\label{principal-results-9}}

\begin{longtable}[]{@{}llll@{}}
\toprule
ID & Result & Status & Review\tabularnewline
\midrule
\endhead
ANS2-L1 & Source formula in extended grammars & PROVED &
V\tabularnewline
ANS2-T02 & Positive drift: depth bounds and abscissa & PROVED &
V\tabularnewline
ANS2-T16 & Growth rate, abscissa, drift dichotomy & PROVED &
V\tabularnewline
ANS2-T05 & Negative drift: abscissa below the characteristic root &
PROVED & V\tabularnewline
ANS2-T13 & Closed countable operator system & PROVED & V\tabularnewline
\bottomrule
\end{longtable}

\hypertarget{recorded-open-questions-8}{%
\subsection{Recorded open questions}\label{recorded-open-questions-8}}

\begin{itemize}
\tightlist
\item
  \textbf{ANS2-O3:} Exact suppression order in the extensions.
\item
  \textbf{ANS2-O5:} Behaviour at the abscissa under negative drift
  (carried on as ANS2N-O2).
\item
  \textbf{ANS2-O6:} Spectrum and continuation of the transfer operators.
\item
  \textbf{ANS2-O7:} Does any extension have a drift exactly zero at
  sigma\_A?
\end{itemize}

\hypertarget{connected-entries-cited-dependencies-9}{%
\subsection{Connected entries (cited
dependencies)}\label{connected-entries-cited-dependencies-9}}

\begin{itemize}
\tightlist
\item
  \textbf{Uses ATL-04} (Collective critical suppression): ANS-T23,
  ANS-T24a, ANS-T27, ANS-T28a.
\item
  \textbf{Used by ATL-11} (Entropy and admissibility): ANS2N-T01,
  ANS2N-T02, ANS2N-T04a, ANS2N-T04b, ANS2N-T05, ANS2N-T06.
\end{itemize}

\hypertarget{review-boundary-9}{%
\subsection{Review boundary}\label{review-boundary-9}}

Every result in this entry has a review verdict.

Qualified verdicts:

\begin{itemize}
\tightlist
\item
  ANS2-T04: Valid for the infimal Gevrey order.
\item
  ANS2-T07: Valid with the precise event definition, now in the text.
\item
  ANS2-T15: Bounded only when the norm is finite (in the text).
\end{itemize}

\hypertarget{all-results-and-proof-sources-9}{%
\subsection{All results and proof
sources}\label{all-results-and-proof-sources-9}}

\begin{itemize}
\tightlist
\item
  \textbf{ANS-II v001:} ANS2-T00, ANS2-L1, ANS2-L2, ANS2-T01, ANS2-T02,
  ANS2-T03, ANS2-T04, ANS2-T05, ANS2-T06, ANS2-T07, ANS2-T08.
\item
  \textbf{ANS-II v002:} ANS2-T14, ANS2-T10, ANS2-T11, ANS2-T12,
  ANS2-T13, ANS2-T15, ANS2-T16, ANS2-T09, ANS2-T09a, ANS2-C16.
\end{itemize}

\hypertarget{atl-11-entropy-and-admissibility}{%
\section{ATL-11: Entropy and
admissibility}\label{atl-11-entropy-and-admissibility}}

\hypertarget{mathematical-object-10}{%
\subsection{Mathematical object}\label{mathematical-object-10}}

Probability laws \(\pi\) on the letters, with entropy \(H(\pi)\) and a
resource functional \(\Lambda(\pi)\) that charges the logarithm of the
multiplier and the depletion of primes. A law is resource-neutral when
no prime drift is negative. The dual variables are prime weights
\(\theta_p\) with optimum \(\theta^\ast\), giving a candidate law
\(\nu^\ast\).

\hypertarget{established-structure-9}{%
\subsection{Established structure}\label{established-structure-9}}

\begin{itemize}
\tightlist
\item
  \textbf{Variational principle.} For \(a>1\), with the suprema taken
  over \textbf{finitely supported} probability laws \(\pi\):
  \(\log\rho=\sup_\pi[H-a\Lambda]\), also over resource-neutral laws,
  and for multipliers above one \(\sigma^\ast=\sup_\pi H/\Lambda\);
  depletion is never optimal (T01, using weights controlled in T02).
  Attainment by a probability law on the complete alphabet is a separate
  result, subject to the criteria below.
\item
  \textbf{Counting.} The three counting functions have the exponent
  \(\sigma^\ast\) for any signs of the drifts (T04a, T04b). Certified
  example (T06).
\item
  \textbf{Optimal law.} Finite optimal laws converge to \(\nu^\ast\);
  criteria for attainment (T03, T07, T08) and a variational inequality
  (T09) lead to the depletion-surplus inequality: \(\nu^\ast\) attains
  if it is neutral, and in particular if all weights are below one half
  (T10).
\item
  \textbf{Shape of the weights.} They vanish identically or are
  supported on infinitely many primes (T05). Under negative drift, ``all
  weights below one half'' and ``weights tend to zero'' cannot both
  hold, in the setting of T11.
\item
  \textbf{At the threshold.} Constraints on large weights (T12a, T12),
  the exceptional decay case (C12), bounds for the moment defect and its
  location at single primes (T13, T14), and a bound on average weight
  (T15).
\end{itemize}

\hypertarget{what-it-reveals-10}{%
\subsection{What it reveals}\label{what-it-reveals-10}}

The exponent of a depleting grammar is an entropy rate under a balance
condition at every prime. Whether the optimum is achieved by a law is a
question about how weight behaves at individual primes.

\hypertarget{further-structure-10}{%
\subsection{Further structure}\label{further-structure-10}}

Failure of attainment, if it occurs, requires a depleted prime with
weight at least one half, on the boundary of the domain (T10) and
visible in the mean divisibility of \(N+1\) at that prime (T14).

\hypertarget{open-boundary-10}{%
\subsection{Open boundary}\label{open-boundary-10}}

Whether the supremum is always attained; convergence and suppression at
the abscissa; lower-order factors of the counts; decay of the weights;
whether a weight can reach one half; and whether the logarithmic average
of the weights vanishes (ANS2N-O1 to O6).

\hypertarget{reading-across-interpretation-10}{%
\subsection{Reading across
(interpretation)}\label{reading-across-interpretation-10}}

This principle is the broadest statement in the collection about
balancing prime resources, and a natural frame for the questions of
ATL-04 and ATL-09. It is a frame only: its hypotheses differ from
theirs. It is also a candidate point of contact with generative
information geometry.

\hypertarget{principal-results-10}{%
\subsection{Principal results}\label{principal-results-10}}

\begin{longtable}[]{@{}llll@{}}
\toprule
ID & Result & Status & Review\tabularnewline
\midrule
\endhead
ANS2N-T01 & Entropy-admissibility variational principle & PROVED &
V\tabularnewline
ANS2N-T03 & Optimal law and attainment criterion & PROVED &
VQ\tabularnewline
ANS2N-T10 & Depletion-surplus inequality & PROVED & V\tabularnewline
ANS2N-T14 & Location of the defect & PROVED & VQ\tabularnewline
\bottomrule
\end{longtable}

\hypertarget{recorded-open-questions-9}{%
\subsection{Recorded open questions}\label{recorded-open-questions-9}}

\begin{itemize}
\tightlist
\item
  \textbf{ANS2N-O1:} Is the entropy supremum always attained?
\item
  \textbf{ANS2N-O2:} Convergence of G at the abscissa, and suppression
  there.
\item
  \textbf{ANS2N-O3:} Lower-order factors of the counting functions.
\item
  \textbf{ANS2N-O4:} Decay of the optimal weights (remaining case:
  maximum exactly 1/2).
\item
  \textbf{ANS2N-O5:} Can an optimal weight reach or exceed 1/2?
\item
  \textbf{ANS2N-O6:} Is the logarithmic average of the optimal weights
  zero?
\end{itemize}

\hypertarget{connected-entries-cited-dependencies-10}{%
\subsection{Connected entries (cited
dependencies)}\label{connected-entries-cited-dependencies-10}}

\begin{itemize}
\tightlist
\item
  \textbf{Uses ATL-04} (Collective critical suppression): ANS-T24a.
\item
  \textbf{Uses ATL-10} (Extended generation and prime-resource balance):
  ANS2-C16, ANS2-L1, ANS2-T16.
\end{itemize}

\hypertarget{review-boundary-10}{%
\subsection{Review boundary}\label{review-boundary-10}}

Every result in this entry has a review verdict.

Qualified verdicts:

\begin{itemize}
\tightlist
\item
  ANS2N-T03: Convergence valid; attainment is open.
\item
  ANS2N-T06: Method valid and stored result reviewed; no independent
  rerun.
\item
  ANS2N-T07: Valid; its hypothesis is superseded by ANS2N-T10.
\item
  ANS2N-T09: Valid; a proof clarification is recorded in ANS-II-ND v003.
\item
  ANS2N-C12: Valid under its hypotheses.
\item
  ANS2N-T14: Proof-writing clarification: in the final Fatou step use a
  limit superior for the drift difference, or a further convergent
  subsequence. The theorem is unaffected; the release file is not
  edited.
\end{itemize}

\hypertarget{all-results-and-proof-sources-10}{%
\subsection{All results and proof
sources}\label{all-results-and-proof-sources-10}}

\begin{itemize}
\tightlist
\item
  \textbf{ANS-II-ND v001:} ANS2N-T02, ANS2N-T01, ANS2N-T03, ANS2N-T04a,
  ANS2N-T04b, ANS2N-T05, ANS2N-T06.
\item
  \textbf{ANS-II-ND v002:} ANS2N-T07, ANS2N-T08, ANS2N-T09, ANS2N-T10,
  ANS2N-T11.
\item
  \textbf{ANS-II-ND v003:} ANS2N-T12a, ANS2N-T12, ANS2N-C12, ANS2N-T13,
  ANS2N-T14, ANS2N-T15.
\end{itemize}

\end{document}
